\documentclass[11pt]{article}
\usepackage[margin=1in]{geometry}
\usepackage{lmodern}
\usepackage[T1]{fontenc}
\usepackage{siunitx}
\usepackage{hyperref}
\setlength{\emergencystretch}{3em}
\begin{document}

\noindent\textbf{Cover Letter — Journal of Chemical Information and Modeling}

\vspace{1em}
\noindent August 10, 2026

\vspace{1em}
Dear Editors,

We submit the manuscript entitled \textit{“Target breadth and mutation resilience in African-natural-product-inspired antimalarial chemotypes: a computational analysis”} for consideration as a research article.

The manuscript presents an integrated computational study connecting African-natural-product-inspired chemical-space exploration with target breadth and mutation-aware resistance hypotheses. Its central methodological premise is that these properties must not be collapsed into a single score: target breadth is represented by a target-specific docking profile, whereas resistance resilience is computed separately from mutation panels using a per-target denominator that excludes non-binding wild-type baselines.

The workflow combines a hybrid library of \num{65856} molecules with a computationally prioritized set of \num{19913} synthesizability-filtered leads and a locked \num{17}-member Set-C cohort. The cohort was evaluated against PfDHFR, PfCRT, PfClpP, and PfATP4 using target-specific structural anchors, yielding \num{68} raw Vina records that passed an automated geometric gate; the underlying poses are archived with full provenance in the repository for independent inspection. Because Vina scores are not calibrated across unlike receptor pockets, the manuscript reports target-wise evidence and does not use a cross-target arithmetic mean. Corrected RRS analysis was derived separately from the declared WT/mutant panel, rather than from the four-target wild-type matrix, and classified the cohort as six A*, five B, five C, and one D profiles. The cross-metric analyses were retained as exploratory and null findings were reported without reinterpretation.

The manuscript addresses a specific prioritization problem: how to combine target breadth and mutation resilience without treating either docking score as experimental activity. Its contribution is a target-specific computational analysis that separates these axes, identifies their joint and discordant profiles, and converts them into testable priorities for biochemical, cellular, mutant-panel, and molecular-dynamics follow-up. The study does not claim measured potency, confirmed target engagement, resistance circumvention, or experimental IC\textsubscript{50}/EC\textsubscript{50}.

Source code, candidate manifests, receptor and ligand files, raw Vina outputs, RRS tables, configuration files, and provenance records are available at \url{https://github.com/NanaEngo/Malaria_codesV2}. All authors have reviewed and approved the manuscript and declare no competing financial interest. This manuscript is not under consideration elsewhere.

Thank you for considering our work.

\vspace{1em}
Sincerely,\\
Serge Guy Nana Engo\\
for all authors

\end{document}
